\chapter{Conclusion} % Chapter 7
\label{Conclusion}

% =============================================================================
% CONCLUSION (3-5 pages)
% Rubric: Clarity of Conclusions (10% of Report)
% =============================================================================

% -----------------------------------------------------------------------------
% SECTION 8.1: Answers to Research Questions
% -----------------------------------------------------------------------------
\section{Answers to Research Questions}

% TODO: Main RQ answer
\subsection{Main Research Question}
%   "How do strategic reasoning, information structure, and enforcement mechanisms
%    each shape the emergence of social order in multi-agent LLM systems?"
%   Direct answer: [to be filled based on results]

% TODO: Sub-question answers
\subsection{Sub-Questions}
%   SQ1 (Reasoning depth × game structure):
%   - [finding — extends Kuusela & Roy to LLM setting]
%
%   SQ2 (Information structure):
%   - [finding — memory/visibility effects]
%
%   SQ3 (Communication & enforcement):
%   - [finding — Leviathan hypothesis]

% -----------------------------------------------------------------------------
% SECTION 8.2: Contributions
% -----------------------------------------------------------------------------
\section{Contributions}
% TODO: Restate contributions with evidence
%   1. Extension of Hobbesian Trap from RL to LLM: [confirmed/modified by results]
%   2. Three mechanism design inputs simultaneously: [demonstrated]
%   3. Per-agent memory as experimental variable: [effect on behavior]
%   4. Phase transitions shift with reasoning depth: [confirmed/not confirmed]
%   5. Reasoning traces as behavioral data: [validation results]
%   6. Methodological: prompt-based reasoning manipulation: [validated]

% -----------------------------------------------------------------------------
% SECTION 8.3: Future Work
% -----------------------------------------------------------------------------
\section{Future Work}

% TODO: Immediate extensions
\subsection{Immediate Extensions}
%   - Complete IV3 (communication & enforcement) if not finished
%   - Full factorial across all three IVs
%   - Scale to larger agent populations (100+)
%   - Test on additional models (cross-architecture validation)

% TODO: Methodological improvements
\subsection{Methodological Improvements}
%   - Full Thought Anchors resampling (compute permitting)
%   - Automated K-level depth measurement
%   - Dynamic game parameters (evolving environment)

% TODO: Broader directions
\subsection{Broader Research Directions}
%   - Real-world multi-agent AI systems (market-making, negotiation)
%   - Hybrid architectures (RL + LLM reasoning)
%   - Connection to AI governance: reasoning depth as a safety-relevant capability
%
% TODO: Koen — DRIELUIK: Van Orde naar Zelfbestuur
%   Gedachte-experiment: autonome AI agents zonder menselijke interventie —
%   welke sociale orde emergeert spontaan?
%   Hobbes' Leviathan als roadmap — van oorlog van allen tegen allen,
%   via het sociaal contract, naar constitutioneel bestuur. Maar bottom-up,
%   zonder Leviathan.
%
%   DEEL 1 — STATE OF NATURE (huidige thesis)
%     Gegeven vaste spelregels, welke sociale structuren ontstaan?
%     Reasoning depth en spelstructuur bepalen of agents in anarchie,
%     coöperatie, of hiërarchie eindigen.
%     Kernbevinding: dieper redeneren transcendeert de spelstructuur niet,
%     maar versterkt het. Cognitieve sophisticatie maakt agents gevoeliger
%     voor hun omgeving, niet onafhankelijker ervan.
%
%   DEEL 2 — SOCIAAL CONTRACT (credible commitment chapter / direct vervolg)
%     Kunnen agents zelf binding mechanisms uitvinden om de Hobbesian trap
%     te ontsnappen? Collateral-gebaseerde contracten als minimale institutie.
%     Test: is er een collateral-drempel waarboven coöperatie stabiel wordt,
%     en verschilt die drempel per reasoning level?
%     Theoretisch: Conitzer (2024) program equilibria, Hobbes's covenant.
%
%   DEEL 3 — CONSTITUTIONELE VORMING (vervolgonderzoek / PhD)
%     Wat als de actieruimte zelf emergent is? Agents componeren acties uit
%     primitieven (Transfer, Condition, Observe, Bind, Scale). Periodiek
%     stemmen ze op innovatiefases waarin nieuwe composities voorgesteld en
%     geadopteerd kunnen worden. De trigger is zelf ook endogeen — agents
%     stemmen over WANNEER er geïnnoveerd wordt.
%
%     Drie condities: vaste acties → composeerbare acties → stembare regels.
%     Meten hoe institutionele agency de emergente orde verandert.
%
%     Drie analyseniveaus:
%     1. Spelen BINNEN regels (Deel 1)
%     2. Regels VERANDEREN (actie-innovatie)
%     3. Beslissen WANNEER regels veranderen (meta-governance)
%
%     Kernvragen:
%     - Welke acties vinden agents uit die de designer niet voorzag?
%     - Blokkeren dominante agents innovatie om de status quo te beschermen?
%     - Bouwen agents instituties die ongelijkheid verminderen, of
%       formaliseren ze bestaande machtsstructuren?
%
%   Theoretische connecties over het drieluik:
%   - Hobbes (Leviathan, 1651): het hele traject van state of nature → contract → constitutie
%   - Ostrom (1990): institutional analysis, rules-in-use, collective choice
%   - Mengesha & Roy (2025): game selection door agents
%   - Hurwicz (1960): mechanism design — agents worden zelf mechanism designers
%   - Het drieluik verbindt alle drie de IVs: reasoning depth bepaalt of
%     agents CAN innoveren, network topology bepaalt of innovaties
%     VERSPREIDEN, en communication scope bepaalt of agents over
%     regelverandering kunnen ONDERHANDELEN.

% -----------------------------------------------------------------------------
% SECTION 8.4: Closing
% -----------------------------------------------------------------------------
\section{Closing Remarks}
% TODO: Brief closing
%   - The "origins of order" are not in the environment alone, nor in the agents alone
%   - They emerge from the interplay of reasoning, information, and institutional structure
%   - As AI agents become more capable and autonomous, understanding these origins
%     becomes critical for designing systems that produce beneficial social outcomes
